\section{Introduction}
\label{sec:intro}

A generated social world needs persistent consequences. In a floating embassy
where endangered languages are political assets, inhabitants must be able to
access archives, trade translation rights, and respond when a flood threatens
the only surviving treaty. These activities require explicit places,
possessions, information, and rules behind the rendered scene.

Several research lines address parts of this problem. Procedural content
generation studies representations and criteria for functional game content
\cite{togelius2011,pcgml}. Word2World and DreamGarden extend language-based
authoring to playable games and development plans \cite{word2world,dreamgarden}.
Generative Agents demonstrates how local memory and encounters can produce
community activities in an authored town \cite{park2023}, while Concordia
grounds language-described actions through a Game Master \cite{vezhnevets2023}.
These advances motivate a joint design question: how can one short premise
specify both a world and the conditions under which a society can act within it?

We present \agora, a system that compiles a one-sentence premise into a
persistent, visually grounded environment for humans and language agents.
Its representation connects rooms, roles, institutions, possessions, private
knowledge, visual assets, and action opportunities. A generated repertoire
makes the world's activities discoverable; an open-proposal channel lets an
actor express a more specific intention. A World Coordinator interprets each
proposal against shared constraints before its effects enter the world state.
For example, proposing a treaty exchange does not by itself transfer ownership:
the resources, participants, and required acceptance must support the change.

The research object is a state trajectory
\begin{equation}
\tau=(S_0,q_1,d_1,S_1,\ldots,q_T,d_T,S_T),
\end{equation}
where $q_t$ is an intended action and $d_t$ is the coordinator decision.
This representation distinguishes an intention, an enacted effect, and a
later response. It also makes explicit which observations concern software
execution, simulated social behavior, or human experience.

Our contributions are threefold:
\begin{enumerate}
\setlength{\itemsep}{2pt}\setlength{\parskip}{0pt}
\item A compositional language-to-world pipeline connecting semantic, visual,
spatial, and mutable social state, with dependency-directed local refinement.
\item A shared action interface for humans and agents that combines generated
opportunities with coordinator-governed proposals and persistent consequences.
\item Audited evidence from ten paired world briefs, a separate 24-cell
one-sentence generation corpus, five locally observed interaction runs, and
21 scripted execution trajectories, distinguishing successful calls from
material consequences and retained failure cases.
\end{enumerate}

The completed studies characterize an authoring and experimentation platform.
A paired memory-intervention design remains prospective
(Appendix~\ref{app:intervention}); behavioral fidelity and causal social effects
require further evidence.

\section{Related Work and Research Position}
\label{sec:related}

\subsection{From Generated Content to Executable Worlds}
Procedural content generation distinguishes what is represented, how it is
generated, and how its fitness is evaluated \cite{togelius2011}. PCGML further
connects learned content models with repair and critique \cite{pcgml}. This
lineage matters for \agora: decomposing a world and repairing a faulty part are
established design ideas. Our use of them ties a repair to explicit dependencies
among role, inventory, geometry, media, and runtime state.

Word2World converts stories into playable tile-based games and evaluates its
generation stages \cite{word2world}. DreamGarden recursively decomposes a user
prompt into plans executed by specialized modules, with user feedback during
Unreal Engine development \cite{dreamgarden}. Both are close precedents for
prompt-driven, staged authoring. \agora emphasizes a reusable social-state
representation: the generated institutions and objects become inputs to
ownership, knowledge access, and action adjudication during play. We make no
claim that prompt-based world generation or hierarchical planning is new.

Environment generation also serves embodied-agent training. ProcTHOR samples
diverse interactive houses \cite{procthor}; Holodeck uses language and spatial
constraints to produce 3D scenes from prompts \cite{holodeck}. Genie instead
learns action-controllable visual dynamics from video \cite{genie}. These
approaches provide distinct substrates for interaction. \agora uses explicit
entity identities and state transitions in a rendered 2D environment, making
social and material changes directly inspectable. Visual generation and
symbolic persistence address complementary aspects of world construction.

\subsection{Grounded Dialogue and Shared Play}
TextWorld provides generated text games with explicit state tracking and
controllable task difficulty \cite{textworld}. ScienceWorld evaluates whether
agents can carry out scientific procedures in an interactive environment
\cite{scienceworld}. Both motivate evaluating executed consequences rather
than the plausibility of a textual plan.

LIGHT is a particularly relevant predecessor: humans and models inhabit a
fantasy environment, combining dialogue with actions grounded in local objects
and characters \cite{light}. Thus shared human--agent play and situated dialogue
are not unique to \agora. Our focus is generating the surrounding social and
material setting from a premise, then coupling visible activities and novel
proposals to the same persistent state. Human-directed events in our logs
establish reachability of the player, while experience and collaboration
quality require a participant study.

\begin{table*}[t]
\centering\small
\setlength{\tabcolsep}{5pt}
\begin{tabularx}{\textwidth}{@{}L{0.17\textwidth}YYY@{}}
\toprule
Research line & Representative substrate & Action or generation mechanism & Relation to \agora \\
\midrule
Word2World; DreamGarden \cite{word2world,dreamgarden} & Generated game content and environments & Story expansion; hierarchical development plans & Precedents for language-driven, staged authoring \\
ProcTHOR; Holodeck \cite{procthor,holodeck} & Procedural or language-conditioned 3D scenes & Interactive objects and spatial constraints & Physical environment diversity and grounding \\
Genie \cite{genie} & Learned visual dynamics & Latent, frame-level control & Complementary representation of world dynamics \\
LIGHT; SOTOPIA \cite{light,zhou2024} & Grounded fantasy play; social scenarios & Dialogue, actions, and social goals & Shared interaction and social evaluation \\
Generative Agents; Concordia \cite{park2023,vezhnevets2023} & Persistent communities; grounded simulations & Memory and planning; Game Master adjudication & Direct precedents for memory and action grounding \\
AgentSociety; LMAgent \cite{agentsociety,lmagent} & Societal and e-commerce environments & Large populations and domain-specific interactions & Scale and behavioral evaluation beyond this study \\
\textbf{\agora} & \textbf{Premise-generated 2D social worlds} & \textbf{Typed shared state; repertoire and open proposals} & \textbf{Connects world authoring to inspectable social trajectories} \\
\bottomrule
\end{tabularx}
\caption{Research positioning by representation and mechanism. Entries describe
the cited systems' emphasis, not an exhaustive feature audit or a performance
ranking. We do not treat unreported capabilities as absent.}
\label{tab:related}
\end{table*}

\subsection{Persistent Communities and Social Simulation}
Generative Agents combines observations, retrieval, reflection, and planning
in a town of 25 agents, illustrating how remembered encounters can support
coordinated activities \cite{park2023}. Concordia provides a flexible simulation
framework in which agents describe intentions and a Game Master grounds their
effects \cite{vezhnevets2023}. The World Coordinator follows this separation of
intention and environment response; \agora connects it to premise-generated
identities, inventories, rules, and rendered spaces.

AgentSociety studies large populations and interventions in a societal
environment \cite{agentsociety}; LMAgent combines multimodal shopping behavior
with large-scale social interaction \cite{lmagent}. Avalon studies cooperation
and confrontation under explicit game rules \cite{avalon}. These systems show
why trajectories and institutions matter beyond isolated dialogue quality.
Our smaller studies address authoring diversity and trace inspectability,
without claiming comparable scale or empirical correspondence to society.

\subsection{Action, Feedback, and Feasibility}
ReAct interleaves reasoning and environment actions \cite{react}; Reflexion
stores verbal feedback for subsequent attempts \cite{reflexion}; Voyager
accumulates executable skills through environment feedback and verification
\cite{voyager}. These methods motivate the observation--intention--outcome
cycle. SayCan explicitly connects language-level usefulness to skill
feasibility \cite{saycan}. \agora applies the analogous distinction to social
and material effects: an intention can be intelligible while its requested
transfer, construction, or status change is infeasible. Its action interface
adds partner acceptance, ownership, and world-specific rules to that decision.
This is a system design connection, not a claim of algorithmic equivalence or
an empirical comparison with those agents.

\subsection{Evaluation and Validity}
SOTOPIA evaluates interaction through multiple social goals and criteria
\cite{zhou2024}. Subsequent work shows that an omniscient simulation can appear
more successful than an information-asymmetric interaction \cite{misleading}.
This motivates our explicit observation boundary and private knowledge, but
does not by itself validate their implementation. SOTOPIA-$\pi$ also reports evaluator overestimation of agents trained for
social interaction \cite{sotopiapi}. Together with broader LLM-judge biases
\cite{judge}, this motivates keeping structural checks, state-derived metrics,
and future human ratings separate.

AgentBoard motivates inspecting intermediate progress rather than relying on
final success alone \cite{agentboard}. We similarly report executed changes
separately from plausible intentions, but do not claim that our mutation
counts measure validated task progress. AI Agents That Matter emphasizes
cost, holdouts, and reproducibility \cite{agentsmatter}; this informs our
separation of unequal-budget architecture comparisons from model coverage
audits and our explicit treatment of retrospective evidence.

The ODD protocol emphasizes explicit entities, scheduling, initialization, and
model rationale for reproducible agent-based modeling \cite{odd}. We adopt
these reporting priorities in Appendix~\ref{app:repro}. Human-grounded agent
simulation, such as models built from interviews and surveys and evaluated
against the people they represent \cite{park1000}, sets a different evidentiary standard from fictional
world coherence. Finally, work on complex contagion distinguishes information
transmission from socially reinforced adoption \cite{centola}. Accordingly,
the prospective intervention measures both warning diffusion and material
treaty rescue; neither network density nor fluent conversation substitutes
for coordinated action.
